% Conclusão (pessoa 6). Convenção igual à do marco teórico 2: sem travessões e sem ponto e vírgula.

\section{Conclusão}
\label{sec:conclusao}

Este trabalho analisou, com base na literatura, como as estimativas por histórias
de usuário sustentam previsões de esforço, prazo e custo em projetos ágeis. A
resposta curta é que elas sustentam, com ressalvas. O \textit{story point} não é
uma unidade monetária e só vira prazo e custo por duas conversões em cadeia. A
\textit{velocity} transforma pontos em tempo, e o custo da equipe transforma
tempo em dinheiro (Seção~\ref{sub:custo}). A Equação~\ref{eq:previsao} é simples
e útil, mas cada um de seus termos herda a incerteza de quem o produziu.

Retomando os objetivos específicos, o marco teórico conceituou histórias de
usuário, \textit{story points} e \textit{velocity} e mostrou como a produtividade
da equipe liga pontos a prazo e a custo. As evidências empíricas reunidas na
Introdução e na Seção~\ref{sec:exemplos} são menos favoráveis do que a prática
sugere. A relação entre pontos e esforço real tem apoio limitado
\cite{tawosi2022}, os modelos de aprendizado de máquina nem sempre superam
referências simples \cite{tawosi2023}, e as avaliações de incerteza tendem ao
excesso de confiança \cite{jorgensen2005}. Os exemplos hipotéticos mostraram
duas aplicações econômicas. Uma faixa de 8 a 10 \textit{sprints} manteve o VPL
positivo, de modo que a incerteza não mudou a decisão de investir, e o custo por
ponto permitiu ordenar funcionalidades pelo retorno, como faz o WSJF.

Daí saem algumas boas práticas para o uso financeiro das estimativas:
%
\begin{alineas}
  \item apresentar o custo como uma faixa, calculada com os extremos da
        \textit{velocity}, e não como um número único (Equação~\ref{eq:intervalo});
  \item dar mais peso ao cenário pessimista, já que a distribuição dos atrasos
        tem cauda longa \cite{little2006};
  \item refazer a previsão a cada \textit{sprint} e acompanhar o custo e o prazo
        com indicadores como $CPI$ e $SPI$ \cite{sulaiman2006};
  \item não comparar pontos nem \textit{velocity} entre equipes e não usar a
        \textit{velocity} como meta de desempenho, o que a faria perder o valor
        preditivo; e
  \item comparar qualquer estimador, humano ou automático, com um
        \textit{baseline} simples, como a mediana histórica \cite{sarro2018}.
\end{alineas}

O trabalho tem limitações. É uma pesquisa bibliográfica e exploratória, com
amostra de fontes desigual entre os periódicos consultados, como já dito nos
procedimentos metodológicos. Os exemplos são hipotéticos e usam valores
escolhidos para ilustrar o cálculo, sem dados de um projeto real. As fontes de
mercado são documentação de fornecedores e não medem a acurácia das previsões.
Esses limites indicam caminhos para estudos futuros, como aplicar o método a
dados históricos de uma organização, para verificar quão larga precisa ser a
faixa de custo, e comparar a previsão por pontos com a previsão em horas ou em
pontos de função em projetos de mesmo tipo.

Em síntese, estimar por histórias de usuário continua sendo uma forma prática e
barata de planejar. O cuidado necessário é não tratar o número que sai da
conta como certo. O que o gestor deve levar à decisão de investimento é uma
faixa de valores, com seus pressupostos explícitos, revista a cada iteração.
